\documentclass[a4paper]{article}
% Uncomment the following line to allow the usage of graphics (.png, .jpg)
\usepackage[pdftex]{graphicx}
% Allow the usage of utf8 characters
\usepackage[utf8]{inputenc}
\usepackage{enumerate}
\usepackage{amssymb}
\usepackage{colortbl}
\usepackage{icomma}
\usepackage[ngerman]{babel}
\usepackage{amsmath}
\usepackage{parskip}
\usepackage{booktabs}
\usepackage{siunitx}
\sisetup{locale=DE} 
\usepackage{geometry}
\geometry{a4paper, top=15mm, left=15mm, right=15mm, bottom=15mm,
	headsep=10mm, footskip=12mm}
\usepackage{href-ul}
\hypersetup{
	colorlinks=true,
	linkcolor=blue,
	urlcolor=blue}
\begin{document}



\title{Humorvolle Physik – Das Erwärmungsgesetz}

	
	\section*{Formel des Tages}
	
	\[
	Q = c \cdot m \cdot \Delta \vartheta
	\]
	
	\begin{itemize}
		\item \( Q \): Wärmeenergie in Joule (J)
		\item \( c \): spezifische Wärmekapazität in \( \frac{\text{J}}{\text{kg} \cdot \text{K}} \)
		\item \( m \): Masse in Kilogramm (kg)
		\item \( \Delta \vartheta \): Temperaturänderung in Kelvin (K) oder Grad Celsius (°C)
	\end{itemize}
	
	\hrulefill
	
	\section*{1. Der faule Kater und die Wärmflasche}
	
Kater Leopold (Masse: 6 kg) liegt auf einer Wärmflasche mit 2 Litern Wasser bei Zimmertemperatur (20 °C). Die Wärmflasche wird auf 60 °C erhitzt, damit seine Pfoten nicht frieren.

	
	\begin{itemize}
		\item a) Wie viel Wärmeenergie wird benötigt, um das Wasser zu erwärmen?  
		\( c_\text{Wasser} = 4{,}18\,\frac{\text{kJ}}{\text{kg}\cdot\text{K}} \)
		
		\vspace{1em}
		\( Q = \underline{\hspace{6cm}} \)
		
		\item b) Bonus: Wenn der Kater 20 \% der Energie schnurrend in Schall umwandelt, wie viel bleibt als Wärme?
		
		\vspace{1em}
		\( Q_\text{Wärme} = \underline{\hspace{6cm}} \)
	\end{itemize}
	
	\hrulefill
	
	\section*{2. Die Spaghetti-Aktion }
	
	Du willst 0{,}5 kg Wasser auf dem Herd zum Kochen bringen – von 15 °C auf 100 °C. Dein Herd hat schlechte Laune und liefert nur 80 \% der eingesetzten Energie als tatsächliche Wärme.
	
	\begin{itemize}
		\item a) Berechne die benötigte Energie bei idealer Umsetzung.
		
		\vspace{1em}
		\( Q_\text{ideal} = \underline{\hspace{6cm}} \)
		
		\item b) Berechne, wie viel Energie du wirklich zuführen musst.
		
		\vspace{1em}
		\( Q_\text{real} = \underline{\hspace{6cm}} \)
	\end{itemize}
	
	\hrulefill
	
	\section*{3. Die Badewanne für Astronauten }
	
	Eine Raumstation simuliert ein gemütliches Bad: 100 kg Wasser werden von 18 °C auf 38 °C erwärmt. Energie ist dort teuer!
	
	\begin{itemize}
		\item a) Berechne die benötigte Wärmemenge.
		
		\vspace{1em}
		\( Q = \underline{\hspace{6cm}} \)
		
		\item b) Wie viele 1000-Watt-Heizstäbe brauchst du, wenn das Bad in 10 Minuten warm sein soll?  
		(1 Watt = 1 Joule pro Sekunde)
		
		\vspace{1em}
		\# Heizstäbe = \underline{\hspace{6cm}}
	\end{itemize}
	
	\hrulefill
	
	\section*{4. Die lauwarme Liebeserklärung }
	
	Ein verliebter Physik-Nerd bringt seiner Angebeteten eine Tasse Tee mit 250 ml Wasser. Der Tee soll auf exakt 50 °C erhitzt werden, kommt aber aus dem Hahn mit 15 °C. Er trägt den Tee 5 Minuten lang durch die Kälte, wo 10 \% der Energie verloren gehen.
	
	\begin{itemize}
		\item a) Berechne die nötige Wärmemenge ohne Verluste.
		
		\vspace{1em}
		\( Q_\text{ideal} = \underline{\hspace{6cm}} \)
		
		\item b) Wie viel Energie muss er tatsächlich zuführen?
		
		\vspace{1em}
		\( Q_\text{real} = \underline{\hspace{6cm}} \)
	\end{itemize}
	

\section*{5. Die heiße Schraube}

Ein Schmied erhitzt eine Eisenschraube mit einer Masse von 300 g von 20 °C auf 180 °C.  
Spezifische Wärmekapazität von Eisen: \( c = 450\, \frac{\text{J}}{\text{kg}\cdot\text{K}} \)

\begin{itemize}
	\item a) Berechne die Wärmeenergie:
	
	\vspace{1em}
	\( Q = \underline{\hspace{6cm}} \)
	
	\item b) Wie groß wäre \( Q \), wenn die Schraube 600 g wiegt?
	
	\vspace{1em}
	\( Q = \underline{\hspace{6cm}} \)
\end{itemize}

\hrulefill

\section*{6. Die schmelzende Butter }
 
 Ein 250 g-Butterblock wird von 6 °C auf 40 °C erwärmt.  
\( c_\text{Butter} = 2100\, \frac{\text{J}}{\text{kg}\cdot\text{K}} \)

\begin{itemize}
	\item a) Berechne die aufgenommene Wärmeenergie:
	
	\vspace{1em}
	\( Q = \underline{\hspace{6cm}} \)
	
	\item b) Welche Temperatur würde 125 g Butter erreichen, wenn man ihr dieselbe Energie zuführt?
	
	\vspace{1em}
	\( \Delta \vartheta = \underline{\hspace{6cm}} \)
\end{itemize}

\hrulefill

\section*{7. Der feurige Ziegelstein }

Ein Ziegelstein (2 kg) wird im Ofen von 25 °C auf 900 °C erhitzt.  
\( c = 840\, \frac{\text{J}}{\text{kg}\cdot\text{K}} \)

\begin{itemize}
	\item a) Berechne die aufgenommene Wärmeenergie:
	
	\vspace{1em}
	\( Q = \underline{\hspace{6cm}} \)
	
	\item b) Zusatzfrage: Welche Rolle spielt die Farbe bei der Aufheizung in der Sonne?
	
	\vspace{1em}
	\underline{\hspace{12cm}}
\end{itemize}

\hrulefill

\section*{8. Der Käsewürfel im Toaster }

Ein Käsewürfel (50 g) wird im Toaster von 10 °C auf 60 °C erwärmt.  
\( c = 3100\, \frac{\text{J}}{\text{kg}\cdot\text{K}} \)

\begin{itemize}
	\item a) Berechne die benötigte Wärmeenergie:
	
	\vspace{1em}
	\( Q = \underline{\hspace{6cm}} \)
	
	\item b) Wie lange braucht ein 150-Watt-Toaster dafür?
	
	\vspace{1em}
	\( t = \underline{\hspace{6cm}} \)
\end{itemize}

\hrulefill

\section*{9. Sauna mit Holzbank }

Ein Holzblock aus Kiefer (1{,}5 kg) erwärmt sich von 25 °C auf 75 °C.  
\( c_\text{Holz} = 1700\, \frac{\text{J}}{\text{kg}\cdot\text{K}} \)

\begin{itemize}
	\item a) Berechne die aufgenommene Wärmeenergie:
	
	\vspace{1em}
	\( Q = \underline{\hspace{6cm}} \)
	
	\item b) Wie viel Energie würde ein gleich schwerer Aluminiumblock aufnehmen?  
	\( c_\text{Alu} = 900\, \frac{\text{J}}{\text{kg}\cdot\text{K}} \)
	
	\vspace{1em}
	\( Q = \underline{\hspace{6cm}} \)
\end{itemize}

\fbox{
	\begin{minipage}{0.5\textwidth}
		Zur Lösung bitte \href{https://www.okuyakl.de/phys/p9qcmL116/ll116.pdf}{hier klicken} oder den QR-Code scannen.\\
		Weitere Arbeitsblätter gibt es unter 
		
		\href{https://www.okuyakl.de}{www.okuyakl.de}
	\end{minipage}
	\hfill
	\begin{minipage}{0.4\textwidth}
		\includegraphics[width=1.5 cm]{../../viecher/zwe03}
		\includegraphics[width=3 cm]{qr116}
		\includegraphics[width=2 cm]{../../viecher/afanticon1}
		
\end{minipage}}

\end{document}

